\documentclass[11pt]{article}
\usepackage[margin=2.4cm]{geometry}
\usepackage{amsmath,amssymb,mathtools}
\usepackage{mathrsfs}
\usepackage{booktabs}
\usepackage[hidelinks]{hyperref}

\newcommand{\pd}{\partial}
\newcommand{\Op}{\operatorname{Op}^{w}}
\newcommand{\Rop}{\mathcal R}
\newcommand{\Top}{\mathcal T}
\newcommand{\Bop}{\mathcal B}
\newcommand{\Lop}{\Lambda}
\newcommand{\eps}{\varepsilon}

\title{A Two-Way Higher-Order Helmholtz Propagator in Orthogonal Terrain-Following Coordinates: A Bremmer--Pseudodifferential Formulation}
\author{Mikhail S. Lytaev}
\date{}

\begin{document}
\maketitle

\begin{abstract}
We formulate a two-way range-marching method for the scalar Helmholtz equation
over irregular terrain. An orthogonal terrain-following map turns the boundary
into a coordinate line and leaves an identity region above the coordinate
transition. A Liouville transformation then writes the equation as
$\varphi_{\xi\xi}+B(\xi)\varphi=0$, where $B$ is a symmetric vertical
pseudodifferential operator. On a propagating spectral subspace, its positive
square root defines flux-normalized forward and backward amplitudes. Their
exact coupled evolution has diagonal one-way generators and a symmetric
Bremmer interaction operator. Iterating the corresponding two-point integral
equations gives successive reflection orders while preserving the distinction
between local higher-order corrections and genuine backscatter. Weyl symbols
identify the principal dispersion relation, the leading coupling, and the
first commutator correction. We then describe a discrete realization that
reuses the metric sampling and high-order vertical rational solves of an
existing terrain-following propagator, while making clear which matrix
fractional-power and two-boundary components a two-way implementation still
requires. This paper gives the formulation and an implementation blueprint;
it does not report numerical results for a completed two-way solver.
\end{abstract}

\section{Introduction}

One-way parabolic-equation (PE) solvers are effective when the field is
predominantly directed along a preferred range coordinate. They omit the
backward-going component, however, and cannot represent reflections generated
by range variation except through errors or ad hoc corrections. Bremmer's
coupling series provides a systematic alternative: decompose the field into
forward and backward waves, propagate each component with a one-way operator,
and couple the two components through an interaction operator
\cite{bremmer,dehoop,gautesen}.

For irregular terrain, the coordinate choice matters to every part of this
construction. Staircase boundaries introduce artificial range discontinuities;
shift maps retain a mixed derivative in the full Helmholtz equation. We use an
orthogonal terrain-following map of the type developed for the high-order PE
in \cite{lytaev_tf}. It straightens the lower boundary, has no mixed
derivative, and becomes the identity above a finite transition layer. This
geometry gives a particularly direct setting for a two-way operator
factorization.

The result below is an analytical and numerical-method formulation, not a
claim that the current JAX terrain-following code already solves the two-way
problem. That code advances one range direction. A two-way solver additionally
needs discrete operator roots, a backward sweep, and a two-point range
boundary solve. Our aim is to specify those pieces in a form that reuses the
existing terrain metric and high-order vertical discretization.

\paragraph{Scope.} We treat a time-harmonic scalar field at fixed frequency,
with constant density in the derivation and a real, smooth wavenumber
$k(\xi,\eta)$. The propagating part of the spectrum is assumed to stay away
from a turning point. Loss, variable density, interfaces, and evanescent
modes can be included in the discrete model, but they remove some of the
self-adjointness and flux-unitarity statements made here. Those extensions
are identified where relevant.

\section{Orthogonal terrain-following Helmholtz equation}

Let $(\xi,\eta)$ denote orthogonal coordinates with scale factors
$h_\xi(\xi,\eta)>0$ and $h_\eta(\xi,\eta)>0$. Write
$r=h_\eta/h_\xi$. In two dimensions,
\begin{equation}\label{eq:laplace}
  \Delta\psi=\frac{1}{h_\xi h_\eta}
  \left[\pd_\xi(r\pd_\xi\psi)+\pd_\eta(r^{-1}\pd_\eta\psi)\right].
\end{equation}
The map is chosen so that $\eta=0$ follows the terrain (or seabed),
$\xi$ labels orthogonal trajectories, and $h_\xi=h_\eta=1$ in an exterior
identity region. The construction in \cite{lytaev_tf} obtains the level curves
from a smoothed harmonic extension of the boundary and follows them with
orthogonal trajectories. In particular, it avoids a mixed $\xi$--$\eta$
derivative in \eqref{eq:laplace}.

The scalar Helmholtz equation $\Delta\psi+k^2\psi=0$ is therefore
\begin{equation}\label{eq:metrichelm}
  \pd_\xi(r\pd_\xi\psi)+\pd_\eta(r^{-1}\pd_\eta\psi)
  +h_\xi h_\eta k^2\psi=0.
\end{equation}
Apply the Liouville scaling
\begin{equation}\label{eq:liouville}
  \psi=r^{-1/2}\varphi.
\end{equation}
After division by $r^{1/2}$, equation \eqref{eq:metrichelm} becomes
\begin{equation}\label{eq:rangeeq}
  \pd_\xi^2\varphi+B(\xi)\varphi=0,
\end{equation}
where the operator in $\eta$ is
\begin{equation}\label{eq:Bdef}
  B\varphi=r^{-1/2}\pd_\eta\!\left[r^{-1}\pd_\eta
  (r^{-1/2}\varphi)\right]+\bigl(k^2h_\xi^2-V\bigr)\varphi,
  \qquad
  V=\frac{\pd_\xi^2(r^{1/2})}{r^{1/2}}.
\end{equation}
For real coefficients and compatible self-adjoint vertical boundary
conditions, $B(\xi)$ is symmetric in $L^2(d\eta)$. The term $V$ is an
explicit range-direction potential. The vertical divergence form in
\eqref{eq:Bdef} must be retained in a discretization to preserve symmetry and
the associated flux identity.

The orthogonal map has an identity exterior. This permits the same physical
outgoing or incoming vertical boundary condition to be imposed above the
terrain transition as in a flat computational domain. It also localizes the
range variation of the metric, a useful property for both the Bremmer
coupling and its numerical truncation.

\section{Weyl symbol and local dispersion}

We use $D_\eta=-i\pd_\eta$ and let $q$ be its dual variable. For a symbol
$a(\xi,\eta,q)$, Weyl quantization is
\begin{equation}\label{eq:weyl}
  \Op(a)u(\eta)=\frac{1}{2\pi}\iint e^{i(\eta-\eta')q}
  a\!\left(\xi,\frac{\eta+\eta'}2,q\right)u(\eta')\,d\eta'\,dq.
\end{equation}
The Weyl product has expansion
\begin{equation}\label{eq:moyal}
 a\# c=ac+\frac{i}{2}\{a,c\}
 -\frac18\left(a_{qq}c_{\eta\eta}-2a_{q\eta}c_{q\eta}
 +a_{\eta\eta}c_{qq}\right)+\cdots,
 \qquad \{a,c\}=a_\eta c_q-a_qc_\eta.
\end{equation}
Thus the symbol of a commutator starts with $i\{a,c\}$, while the first
derivative term cancels for the Weyl symbol of a square.

For $D=-i\pd_\eta$ and smooth $g$, the identity
\begin{equation}\label{eq:DgD}
  DgD=\Op\!\left(gq^2+\frac14g_{\eta\eta}\right)
\end{equation}
gives the exact Weyl symbol of \eqref{eq:Bdef}:
\begin{equation}\label{eq:bexact}
  B=\Op(b),\qquad
  b=k^2h_\xi^2-\frac{q^2}{r^2}-V-
  \frac{r_\eta^2}{4r^4}.
\end{equation}
The principal symbol is
\begin{equation}\label{eq:bprincipal}
  b_0=k^2h_\xi^2-\frac{q^2}{r^2}.
\end{equation}
The characteristic equation of \eqref{eq:rangeeq} is
$\kappa^2=b_0$, with $\kappa$ dual to $\xi$. In terms of the physical
vertical wavenumber $q/h_\eta$ and range wavenumber $\kappa/h_\xi$, this is
\begin{equation}\label{eq:dispersion}
  \left(\frac{\kappa}{h_\xi}\right)^2+
  \left(\frac{q}{h_\eta}\right)^2=k^2.
\end{equation}
For a ray angle $\theta$ measured from the level curve,
$q=kh_\eta\sin\theta$ and
$\sqrt{b_0}=kh_\xi\cos\theta$. The forward and backward characteristic
sheets are separated by $2\sqrt{b_0}$ wherever the propagating aperture
stays away from grazing incidence.

If metric coefficients vary on length $\ell$ and the characteristic
wavenumber is $k$, the semiclassical parameter is $\eps=(k\ell)^{-1}$.
The two zeroth-order corrections in \eqref{eq:bexact}, $V$ and
$r_\eta^2/(4r^4)$, are of relative order $\eps^2$ for a smoothly varying
metric. The exact symbol is real; it has no order-$\eps$ imaginary
subprincipal part in this Liouville gauge.

\section{Exact forward--backward splitting}

On a spectral subspace where $B(\xi)$ is positive and bounded away from zero,
define its positive roots
\begin{equation}\label{eq:roots}
  \Lambda=B^{1/2},\qquad Q=B^{1/4},\qquad Q^2=\Lambda.
\end{equation}
All powers here are operator powers at the same $\xi$. Decompose the field
and its range derivative as
\begin{equation}\label{eq:split}
  \varphi=Q^{-1}(u_++u_-),\qquad
  \varphi_\xi=iQ(u_+-u_-),
\end{equation}
or, equivalently,
\begin{equation}\label{eq:splitinv}
  u_\pm=\frac12\left(Q\varphi\mp iQ^{-1}\varphi_\xi\right).
\end{equation}
The signs are fixed so that for a constant $B$, $u_+$ evolves as
$e^{+i\Lambda\xi}$ and $u_-$ as $e^{-i\Lambda\xi}$.

Differentiate \eqref{eq:splitinv}, use \eqref{eq:rangeeq}, and retain the
$\xi$-dependence of $Q$. This gives the exact coupled system
\begin{equation}\label{eq:coupled}
  \pd_\xi\begin{pmatrix}u_+\\u_-\end{pmatrix}=
  \begin{pmatrix}i\Lambda+T & \Rop\\
                  \Rop & -i\Lambda+T\end{pmatrix}
  \begin{pmatrix}u_+\\u_-\end{pmatrix},
\end{equation}
with
\begin{equation}\label{eq:TRexact}
  T=\frac12[Q_\xi,Q^{-1}],\qquad
  \Rop=\frac12(Q_\xi Q^{-1}+Q^{-1}Q_\xi).
\end{equation}
For real self-adjoint $Q$, $T^*=-T$ and $\Rop^*=\Rop$. Hence
\begin{equation}\label{eq:flux}
  \pd_\xi\left(\|u_+\|^2-\|u_-\|^2\right)=0.
\end{equation}
The invariant is the net range-directed flux. The individual squared norms
are the forward and backward fluxes in this normalization.

Equation \eqref{eq:coupled} is exact on the chosen positive spectral subspace.
The one-way approximation sets $u_-=0$ and discards the feedback through
$\Rop$. The two-way model retains it. If $B$ has evanescent or nonpositive
eigenvalues, the positive-root construction and the unitary flux
interpretation do not apply to those modes. A practical implementation must
either separate them from the propagating aperture, use a complex branch with
appropriate incoming/outgoing boundary conditions, or work with an alternative
first-order state-space factorization.

\section{Pseudodifferential expansion of the coupling}

Let $\lambda=b^{1/2}$ on the propagating symbol region. Functional calculus
gives, to principal Weyl order,
\begin{equation}\label{eq:symbolroots}
  \sigma(\Lambda)=\lambda+O(\eps^2),\qquad
  \sigma(Q)=b^{1/4}+O(\eps^2).
\end{equation}
The first correction to the square-root symbol follows directly from
$\lambda_W\#\lambda_W=b$. With $\lambda_0=\sqrt b$,
\begin{equation}\label{eq:sqrtcorrection}
  \sigma(\Lambda)=\lambda_0+\lambda_2+O(\eps^4),\qquad
  \lambda_2=\frac{(\lambda_0)_{\eta\eta}(\lambda_0)_{qq}
  -(\lambda_0)_{\eta q}^{,2}}{8\lambda_0}.
\end{equation}
This term records the phase-space curvature of the dispersion surface. It
can be included in the diagonal phase when the leading-symbol rational
approximation is not accurate enough; equivalently, one can construct a
corrected symbol of $L_\beta$ from the same operator functional calculus.
The anticommutator in \eqref{eq:TRexact} has the principal symbol
\begin{equation}\label{eq:Rsymbol}
  \rho:=\sigma(\Rop)=\frac{b_\xi}{4b}+O(\eps^3),
\end{equation}
where the displayed term is order $\eps$ relative to $\lambda$. The first
Poisson-bracket contribution cancels in this anticommutator. For the
commutator,
\begin{equation}\label{eq:Tsymbol}
  \sigma(T)=\frac{i}{32}\frac{\{b,b_\xi\}}{b^2}+O(\eps^3),
\end{equation}
which begins at order $\eps^2$ relative to the principal one-way rate. The
commutator is skew-adjoint and contributes a phase/transport correction to
each directional component.

At leading order the interaction is therefore the logarithmic range
variation of the local longitudinal wavenumber:
\begin{equation}\label{eq:rhoangle}
  \rho\simeq\frac14\pd_\xi\log b_0.
\end{equation}
For a narrow ray packet this reduces to the corresponding derivative of
$kh_\xi\cos\theta$. The coupling vanishes in the identity exterior when the
medium is also range-independent there.

\paragraph{Higher-order directional diagonalization.} A higher-order
one-way symbol can be constructed by eliminating the off-diagonal coupling
order by order. To state the recursion, write
$A_+=i\Lambda+T$ and $A_-=-i\Lambda+T$ and represent a locally forward
subspace as a graph $u_-=K(\xi)u_+$. Invariance of this graph under
\eqref{eq:coupled} requires the operator Riccati equation
\begin{equation}\label{eq:graphriccati}
  K_\xi=\Rop+A_-K-KA_+-K\Rop K.
\end{equation}
If $k$ is the Weyl symbol of $K$, this becomes the symbol Riccati equation
\begin{equation}\label{eq:graphriccatisymbol}
  k_\xi=\rho+a_-\# k-k\# a_+-k\#\rho\# k,
  \qquad a_\pm=\sigma(A_\pm).
\end{equation}
The leading symbol equation has denominator $2i\lambda$; it is solvable as
long as the forward and backward sheets remain separated. Expanding $K$ in
powers of $\eps$ and substituting back into the $u_+$ equation produces
successive local corrections to a one-way generator. This is the
pseudodifferential counterpart of higher-order WKB factorization
\cite{fishman,stolk}. It does not replace the nonlocal two-way solution near
a turning point or a strong reflector.

\section{Bremmer series and boundary-value formulation}

Let $G_+(\xi,s)$ and $G_-(\xi,s)$ be evolution operators for the diagonal
generators $A_+$ and $A_-$:
\begin{equation}\label{eq:G}
  \pd_\xi G_\pm(\xi,s)=A_\pm(\xi)G_\pm(\xi,s),\qquad
  G_\pm(s,s)=I.
\end{equation}
For real self-adjoint $B$, these propagators are unitary on the propagating
subspace. Given an incident forward field at $\xi=0$ and no incoming
backward field at $\xi=X$, the exact variation-of-constants equations are
\begin{align}
  u_+(\xi)&=G_+(\xi,0)u_+(0)+
  \int_0^\xi G_+(\xi,s)\Rop(s)u_-(s)\,ds,\label{eq:volterra+}\\
  u_-(\xi)&=-\int_\xi^X G_-(\xi,s)\Rop(s)u_+(s)\,ds.
  \label{eq:volterra-}
\end{align}
These equations form a two-point boundary-value problem. Substitution of
\eqref{eq:volterra-} into \eqref{eq:volterra+}, followed by iteration, yields
the Bremmer series. Each additional round trip introduces two interactions
with $\Rop$ and one additional pair of directional propagations. A first
backward approximation follows by substituting the zeroth forward field
$u_+^{(0)}(s)=G_+(s,0)u_+(0)$ into \eqref{eq:volterra-}:
\begin{equation}\label{eq:firstreflection}
  u_-^{(1)}(\xi)=-\int_\xi^X G_-(\xi,s)\Rop(s)G_+(s,0)u_+(0)\,ds.
\end{equation}

For a single propagating mode with phase $\Phi_\xi=\lambda>0$, the first
reflection has the oscillatory form
\begin{equation}\label{eq:oscillatory}
  u_-^{(1)}(\xi)=-e^{-i\Phi(\xi)}
  \int_\xi^X \rho(s)e^{2i\Phi(s)}u_+(0)\,ds.
\end{equation}
Integration by parts separates a locally slaved component, of relative size
$O(\eps)$, from a rapidly oscillating reflected remainder. Feeding the local
component back into the forward equation produces order-$\eps^2$ corrections,
including the familiar second-order WKB terms. Genuine backscatter is
controlled by the regularity of the range variation and the separation of
the two characteristic sheets. With $C^m$ coefficients, compactly supported
range variation, and no stationary phase point, repeated integration by parts
gives algebraic decay with the available number of derivatives. A
$C^\infty$ profile gives superalgebraic decay; an exponential estimate
requires analytic continuation and corresponding contour-deformation
hypotheses. A merely $C^2$ terrain map does not justify an
$O(\eps^\infty)$ or exponential claim.

For the operator series, a sufficient small-coupling condition is
\begin{equation}\label{eq:smallcoupling}
  \int_0^X\|\Rop(s)\|\,ds<1.
\end{equation}
Under unitary diagonal propagation, this bounds successive integral
iterations. It is sufficient rather than necessary; near grazing, a turning
point, a sharp interface, or a strong reflector the series may require a
different preconditioner or a direct block boundary solve.

\section{Weyl-symbol propagators and rational approximation}

\subsection{Symbol evolution and Bremmer terms}

Let $\mathbf u=(u_+,u_-)^T$ and let $\mathbf a$ be the matrix Weyl symbol of
the generator in \eqref{eq:coupled}:
\begin{equation}\label{eq:amatrix}
  \mathbf a=\begin{pmatrix}a_+&\rho\\\rho&a_-\end{pmatrix},
  \qquad a_+=i\lambda+t,\quad a_-=-i\lambda+t,
\end{equation}
where $\lambda=\sigma(\Lambda)$, $t=\sigma(T)$, and $\rho=\sigma(\Rop)$.
For a $2\times2$ symbol matrix, $\# $ denotes matrix multiplication with
Weyl products in the entries. The propagator symbol $\mathbf g(\xi,s)$ is
defined by
\begin{equation}\label{eq:starprop}
  \pd_\xi\mathbf g(\xi,s)=\mathbf a(\xi)\#\mathbf g(\xi,s),
  \qquad \mathbf g(s,s)=\mathbf I.
\end{equation}
Thus the exact symbol solution is a range-ordered \emph{star exponential},
not the pointwise exponential of the integrated symbol. At leading
semiclassical order, $\#$ becomes ordinary multiplication and
\begin{equation}\label{eq:leadingprop}
  \pd_\xi\mathbf g_0=\mathbf a_0\mathbf g_0.
\end{equation}
The next terms follow by substituting the Moyal expansion \eqref{eq:moyal}
into \eqref{eq:starprop}; they include phase-space derivatives and account
for noncommutation of the vertical operators at different ranges.

The diagonal propagator symbols $g_\pm(\xi,s)$ satisfy
\begin{equation}\label{eq:diagonalstar}
  \pd_\xi g_\pm(\xi,s)=a_\pm(\xi)\# g_\pm(\xi,s),
  \qquad g_\pm(s,s)=1.
\end{equation}
The first reflected field is the operator with symbol
\begin{equation}\label{eq:firstreflectionsymbol}
  s_-^{(1)}(\xi)=-\int_\xi^X
  g_-(\xi,s)\#\rho(s)\#g_+(s,0)\,ds,
  \qquad u_-^{(1)}(\xi)=\Op(s_-^{(1)}(\xi))u_+(0).
\end{equation}
Higher Bremmer orders are nested range integrals of alternating $g_+$ and
$g_-$ symbols, with a factor $\rho$ at every change of direction and with
all products interpreted as $\#$. For a single isolated ray, the leading
star products reduce to scalar products and \eqref{eq:firstreflectionsymbol}
reduces to the oscillatory reflection integral \eqref{eq:oscillatory}.

\subsection{Frozen-range propagator}

For a short step $d=\Delta\xi$, first freeze the coefficients at the
midpoint $\xi_{j+1/2}$. The diagonal operator generators are
$A_\pm=\pm i\Lambda+T$. A symmetric splitting separates the skew-adjoint
commutator correction from the principal phase:
\begin{equation}\label{eq:splitstep}
  G_\pm(\xi_{j+1},\xi_j)\simeq
  e^{dT_{j+1/2}/2}\,e^{\pm id\Lambda_{j+1/2}}
  e^{dT_{j+1/2}/2}+O(d^3),
\end{equation}
for frozen smooth coefficients. A higher-order Magnus or symmetric
composition can be used when the range variation over one step is not
negligible. The middle factor is a function of the vertical operator. Choose
a reference wavenumber $\beta>0$ and define
\begin{equation}\label{eq:Lop}
  L_\beta=\frac{B}{\beta^2}-I,
  \qquad \ell_\beta(\xi,\eta,q)=\frac{b(\xi,\eta,q)}{\beta^2}-1.
\end{equation}
Then
\begin{equation}\label{eq:phasefunction}
  e^{\pm id\Lambda}=e^{\pm i\beta d}
  \exp\!\left[\pm i\beta d\left(\sqrt{I+L_\beta}-I\right)\right].
\end{equation}
This is the same reduced square-root phase used in Lytaev's JCS22 one-way
propagator \cite{lytaev_jcs}, with $L_\beta$ now the Liouville-transformed
terrain-following vertical operator.

\subsection{Rational approximation in Weyl form}

Let $\mathscr R_\pm(\zeta)$ be a rational approximation on the propagating
spectral interval $\mathcal I$ to
\begin{equation}\label{eq:targetscalar}
  F_\pm(\zeta)=
  \exp\!\left[\pm i\beta d\left(\sqrt{1+\zeta}-1\right)\right],
  \qquad \zeta\in\mathcal I.
\end{equation}
For example, in factored form,
\begin{equation}\label{eq:rational}
  \mathscr R_\pm(\zeta)=
  \prod_{m=1}^{p}\frac{1+a_m^\pm\zeta}{1+b_m^\pm\zeta},
  \qquad
  \sup_{\zeta\in\mathcal I}|F_\pm(\zeta)-\mathscr R_\pm(\zeta)|\le\delta.
\end{equation}
The coefficients may come from a local Pad\'e fit about $\zeta=0$ or from
an interval fit such as Chebyshev--Pad\'e or rational interpolation. The
interval is selected from the desired angle and wavenumber aperture, as in
JCS22. In terrain-following coordinates its endpoints must include the
metric-scaled spectrum $b_0\simeq k^2h_\xi^2-q^2/r^2$ (and the intended
subprincipal corrections if they are retained).

At the operator level, set
\begin{equation}\label{eq:rationaloperator}
  \mathscr R_\pm(L_\beta)=
  \prod_{m=1}^{p}(I+a_m^\pm L_\beta)
  \left[\prod_{m=1}^{p}(I+b_m^\pm L_\beta)\right]^{-1}.
\end{equation}
Each factor requires one shifted solve. For example,
$(I+b_m^\pm L_\beta)v_m=(I+a_m^\pm L_\beta)w_{m-1}$.
With the flux-form finite-difference realization of $B$, these are the
tridiagonal vertical solves already used by the high-order terrain-following
propagator. The two directions use conjugate rational fits for real $B$,
so the discrete frozen-range phase pair has the expected adjoint relation.

The exact Weyl symbol of a rational function is most cleanly written using
star products. For a polynomial $p(\zeta)=\sum_j p_j\zeta^j$, define
$p^{\#}(\ell_\beta)=\sum_jp_j\ell_\beta^{\#j}$, with
$\ell_\beta^{\#0}=1$ and $\ell_\beta^{\#(j+1)}=
\ell_\beta^{\#j}\#\ell_\beta$. If $p_\pm(\zeta)$ and $q_\pm(\zeta)$ are the
numerator and denominator polynomials in \eqref{eq:rational}, form their
star-functional-calculus symbols $p_\pm^{\#}(\ell_\beta)$ and
$q_\pm^{\#}(\ell_\beta)$.
Then
\begin{equation}\label{eq:rationalsymbol}
  \sigma^w\!\left(\mathscr R_\pm(L_\beta)\right)=
  p_\pm^{\#}(\ell_\beta)\#
  \left(q_\pm^{\#}(\ell_\beta)\right)_{\#}^{-1},
  \qquad
  q_\pm^{\#}(\ell_\beta)\#
  \left(q_\pm^{\#}(\ell_\beta)\right)_{\#}^{-1}=1.
\end{equation}
The star inverse exists on the chosen symbol region when the rational
denominator has no zero there. Since the first Weyl correction to a scalar
function of one real symbol vanishes, the principal symbol is simply
\begin{equation}\label{eq:rationalprincipal}
  \sigma^w\!\left(\mathscr R_\pm(L_\beta)\right)
  =\mathscr R_\pm(\ell_\beta)+O(\eps^2).
\end{equation}
For the commutator half-step write
$\tau_d=\sigma^w(e^{dT/2})$; at principal order
$\tau_d=e^{dt/2}+O(\eps^2)$.
Combining \eqref{eq:splitstep}--\eqref{eq:rationalprincipal} gives the
frozen-range Weyl symbol of the rationalized directional propagator:
\begin{equation}\label{eq:Gsymbolrational}
  g_{\pm,j}^{\rm rat}=
  \tau_d\#\left[e^{\pm i\beta d}
  p_\pm^{\#}(\ell_\beta)\#
  \left(q_\pm^{\#}(\ell_\beta)\right)_{\#}^{-1}\right]
  \#\tau_d+O(d^3),
\end{equation}
where all symbols are evaluated at $\xi_{j+1/2}$. At principal order,
\begin{equation}\label{eq:Gsymbolerror}
  g_{\pm,j}^{\rm rat}\simeq
  \exp\!\left[\pm id\sqrt{b(\xi_{j+1/2},\eta,q)}\right],
  \qquad
  |\text{phase error}|\le\delta+O(\eps^2)+O(d^3).
\end{equation}
The local Pad\'e fit concentrates accuracy near normal incidence
($\ell_\beta=0$); an interval fit controls the error over the full symbol
support. This symbol-level error is distinct from the vertical finite
difference dispersion error and from the Bremmer-series truncation error.

In flat Cartesian coordinates $h_\xi=h_\eta=r=1$, let the local wavenumber
be $k_0n(\xi,\eta)$ and choose reference $\beta=k_0$. Then
$b=k_0^2n^2-q^2$ and
\begin{equation}\label{eq:jcsrecovery}
  \ell_{k_0}=-\frac{q^2}{k_0^2}+(n^2-1).
\end{equation}
For $n=1$ this reduces to $-q^2/k_0^2$, exactly the homogeneous-medium case;
for variable $n$ it is the scalar argument used in JCS22. Equation
\eqref{eq:rationaloperator} is therefore its
terrain-following operator generalization, while
\eqref{eq:rationalsymbol}--\eqref{eq:Gsymbolerror} give the corresponding
Weyl-symbol derivation.

\section{Discrete realization in terrain-following coordinates}

\subsection{Vertical operator and one-way building blocks}

On a range grid $\xi_j=j\Delta\xi$ and vertical grid $\eta_n=n\Delta\eta$,
construct a matrix $B_j$ from the flux-form operator \eqref{eq:Bdef}. The
discrete vertical scheme should preserve the self-adjoint weighting and
boundary conditions. In the existing terrain-following JAX propagator, the
coefficient arrays include the metric weight $m=r^2$, the flux coefficient
$r$ (or its density-weighted version), the local wavenumber contrast, and
midpoint flux coefficients. Each rational substep is a tridiagonal solve;
the range step also includes a flux-normalizing fourth-root pair
\cite{lytaev_tf}.

Those pieces can be reused to build actions of $B_j$, and rational resolvent
solves can approximate $B_j^{1/2}$ and $B_j^{1/4}$ on a specified propagating
spectrum. Alternatively, for small verification problems, a symmetric
matrix eigendecomposition gives a reference implementation. Define
\begin{equation}\label{eq:discretecoupling}
  Q_j=B_j^{1/4},\qquad \Lambda_j=B_j^{1/2},\qquad
  \Rop_{j+1/2}=\frac12\left(
  \frac{Q_{j+1}-Q_j}{\Delta\xi}Q_{j+1/2}^{-1}
  +Q_{j+1/2}^{-1}\frac{Q_{j+1}-Q_j}{\Delta\xi}\right),
\end{equation}
with a consistent midpoint or higher-order approximation for $Q_{j+1/2}$.
The finite difference in \eqref{eq:discretecoupling} approximates $Q_\xi$;
for high accuracy it should be replaced by a symmetric multi-column
differentiation or a direct discrete derivative of the matrix-function
action. A discontinuous staircase metric is unsuitable for this derivative
because it creates artificial impulses in the coupling.

\subsection{Two directional sweeps}

Let $\mathcal{P}_+$ and $\mathcal{P}_-$ be discrete one-way propagators for
$A_+$ and $A_-$, including the skew-adjoint term $T$ to the desired order.
Given a current forward field, a backward sweep evaluates the discrete
version of \eqref{eq:volterra-}; the next forward sweep evaluates
\eqref{eq:volterra+}. Starting with $u_-^{(0)}=0$ gives the successive
Bremmer reflection corrections. An equivalent approach solves the block
tridiagonal discretization of \eqref{eq:coupled} directly, which avoids
fixed-point convergence issues and provides a reference for truncated-series
results.

\begin{center}
\begin{tabular}{@{}lll@{}}
\toprule
Component & Existing terrain-following code & Required two-way extension\\
\midrule
Map and metric & Orthogonal map, $h_\xi,h_\eta,r$ & Reuse unchanged\\
Vertical operator & High-order flux-form rational solves & Expose symmetric $B_j$ action\\
Range step & Forward rational PE step & Separate $G_+$ and $G_-$\\
Flux scaling & Approximate $B^{\pm1/4}$ column pair & Root actions in the split state\\
Backscatter & Not represented & $\Rop$ and backward range sweep\\
Range boundary & One-way source/outgoing top & Incident data at $0$, no incoming at $X$\\
\bottomrule
\end{tabular}
\end{center}

\subsection{Boundary conditions and absorbing regions}

The terrain boundary condition is imposed on the coordinate line $\eta=0$.
For a finite-impedance boundary, the physical normal derivative must be
transformed with the local metric; the current solver's compact closure is a
candidate discretization for this part. At the top, the orthogonal map is the
identity, so a vertical outgoing condition can be imposed without map
corrections. For the range boundary, one prescribes the incident $u_+(0)$ and
sets the incoming-from-the-right component $u_-(X)$ to zero. This range
condition is distinct from the vertical top boundary condition.

An absorbing layer makes the discrete operator non-self-adjoint. Then the
flux conservation and unitary bounds above no longer hold exactly. A robust
implementation should first validate the two-way solver on a lossless finite
domain with exact outgoing vertical conditions. Absorption can then be added
with a monitored reflection coefficient and a biorthogonal or energy estimate
appropriate to the damped operator.

\section{Verification programme}

A two-way implementation should be developed in stages.
\begin{enumerate}
\item \textbf{Constant flat medium.} Check that $\Rop=0$, the two components
  propagate with the expected signs, and their superposition agrees with the
  exact Helmholtz solution.
\item \textbf{Range-dependent one-dimensional medium.} With no vertical
  variation, compare the full two-way result and successive Bremmer orders
  against an exact transfer-matrix or boundary-value solution. Measure both
  the reflected amplitude and convergence with range step.
\item \textbf{Smooth terrain map.} Use a smooth isolated hill and compare
  zero-, first-, and higher-reflection results with a direct two-way
  discretization. Refine the map and grid independently to separate terrain
  regularity error from propagation error.
\item \textbf{Impedance and evanescence.} Add finite ground impedance, then
  evanescent modes and absorption, checking flux balance where applicable
  and boundary reflection where flux conservation does not apply.
\end{enumerate}
No numerical accuracy or speed claim is made here; the existing one-way JAX
solver is a source of reusable discretization components, not a two-way
reference solution.

\section{Conclusion}

Orthogonal terrain-following coordinates put the terrain-boundary Helmholtz
equation in a form suited to directional operator factorization. The
Liouville-transformed equation has a symmetric vertical operator whose
fractional powers define flux-normalized forward and backward components.
Their exact coupled evolution yields a Bremmer series, while Weyl calculus
identifies the leading interaction and its higher-order corrections.

The existing high-order terrain-following PE contributes the coordinate map,
metric and flux-form vertical solves. A genuine two-way higher-order
propagator still requires discrete square-root actions, consistent coupling
operators, and a two-point range boundary solution. Deriving and validating
those components is the next step; reversing the existing one-way march alone
does not produce a Bremmer solver.

\begin{thebibliography}{10}
\bibitem{bremmer}
H. Bremmer, ``The W.K.B. approximation as the first term of a geometric-optical
series,'' \emph{Communications on Pure and Applied Mathematics}, vol. 4,
pp. 105--115, 1951.

\bibitem{dehoop}
M. V. de Hoop, ``Generalization of the Bremmer coupling series,''
\emph{Journal of Mathematical Physics}, vol. 37, no. 7, pp. 3246--3282, 1996,
doi: 10.1063/1.531566.

\bibitem{gautesen}
A. K. Gautesen and M. V. de Hoop, ``Uniform asymptotic expansion of the
generalized Bremmer series,'' \emph{SIAM Journal on Applied Mathematics},
vol. 60, no. 4, pp. 1302--1329, 2000, doi: 10.1137/S0036139997323619.

\bibitem{fishman}
L. Fishman and J. J. McCoy, ``Derivation and application of extended
parabolic wave theories. I. The factorized Helmholtz equation,''
\emph{Journal of Mathematical Physics}, vol. 25, no. 2, pp. 285--296, 1984.

\bibitem{stolk}
C. C. Stolk, ``A pseudodifferential equation with damping for one-way wave
propagation in inhomogeneous acoustic media,'' \emph{Wave Motion}, vol. 40,
no. 2, pp. 111--121, 2004, doi: 10.1016/j.wavemoti.2003.12.016.

\bibitem{lytaev_jcs}
M. S. Lytaev, ``Rational interpolation of the one-way Helmholtz propagator,''
\emph{Journal of Computational Science}, vol. 58, article 101536, 2022,
doi: 10.1016/j.jocs.2021.101536.

\bibitem{zworski}
M. Zworski, \emph{Semiclassical Analysis}. Providence, RI: American
Mathematical Society, 2012.

\bibitem{lytaev_tf}
M. S. Lytaev, ``Orthogonal terrain-following coordinates for the higher-order
parabolic equation: Wide-angle propagation over irregular terrain with
discrete transparent boundaries,'' manuscript accompanying the
terrain-following implementation, \texttt{pywaveprop/helmholtz\_tf\_jax.py}.
\end{thebibliography}

\end{document}
